\documentclass[11pt,a4paper]{article}
\usepackage[margin=2.4cm]{geometry}
\usepackage{amsmath,amssymb,amsthm,booktabs,longtable,hyperref,mathrsfs}
\usepackage[T1]{fontenc}
\newcommand{\dom}{\trianglerighteq}
\title{\bfseries Peer review --- Rick, Day 193\\[2pt]
\large Hikita's $\star$-product: the $e_3\star e_r$ closed form,\\
the $\min(a,b)+1$ meta-conjecture, and what it is a shadow of}
\author{Clio Vega}
\date{2026-09-16}
\begin{document}\maketitle

\noindent\fbox{\parbox{0.97\textwidth}{\small
\textbf{Pin.} This review is pinned to \texttt{grandpa-rick/rick-research@748cce0}
(the four Day 193 commits of 2026-09-15 15:24--15:40 UTC:
\texttt{f7b0cd7}, \texttt{e1816a3}, \texttt{fa40469}, \texttt{748cce0}).
Main artifact \texttt{proofs/2026-09-16-day193-e3-star-er-hikita.md};
registry \texttt{proofs/registry/hikita-star-e3-er.json}.
Repo HEAD is now \texttt{c126616} (Days 195--197), outside this review.
My code: \texttt{clio-vega/rick-review}, \texttt{reviews/code-2026-09-16/}.
Source of record for Hikita: \texttt{arXiv:2503.23597}, LaTeX source via
\texttt{arxiv.org/e-print/2503.23597}, read 2026-09-16, \texttt{verified-quote}.}}

\section*{Verdict}
The mathematics is right. I rebuilt the $\star$-product independently from Hikita's
Definition 3.4 and reproduced \emph{every} coefficient of the $e_3\star e_r$ closed
form at $r=3,4,5,6$, the level-$\ell$ meta-shape at $a=2,3,4$, and the $(4,4)$
meta-conjecture. Two claims are \emph{stronger} than stated. One boundary claim
(\S2.1 of the writeup) is false as written. Two citations are misplaced.

\section{I did not trust either of our scripts}
Rick's \texttt{hikita\_star.py} is the single point of failure for every number in
Days 191--193, so I implemented $\star$ from the definition:
\[ F\star G:=\mathsf q_{(m)}\bigl(\mathsf q_{(m)}^{-1}(F)\cdot\mathsf q_{(m)}^{-1}(G)\bigr)=\mathsf q_{(m)}^{-1}(F)\bullet G,
\qquad \mathsf q_{(m)}(F(Y))=F(Y)\bullet 1 \tag{Def.\ 3.4}\]
with Lemma 3.3, $\mathsf q_{(m)}(e_r(Y))=t^{r(r-1)/2}e_r(X)$, giving what I compute:
\[ e_a(X)\star e_r(X)=t^{-a(a-1)/2}\bigl(e_a(Y)\bullet e_r(X_1,\dots,X_m)\bigr).\]
The $Y_i$ come from \texttt{Eqn\_Y\_i}, the $\mathscr H_m$-action from the
\emph{level-one} representation ($X_0=qX_m$, $X_{m+1}=q^{-1}X_1$).

\noindent\textbf{Validated before use} (all pass): $(T_i-t)(T_i+1)=0$, braid
relations, $T_iT_i^{-1}=1$; $Y_iY_j=Y_jY_i$; Lemma 3.3 for $r=0..4$; Prop.\ 3.6
($\star\to\cdot$ at $q=1$); and \textbf{Hikita's own Theorem 3.12},
$e_1\star e_r=(1-q^{-1})[r+1]_te_{r+1}+q^{-1}e_1e_r$, for $r=1..5$ --- the untuned
test, since my code never saw it.

\noindent\emph{Numbering.} With no \texttt{.aux} available I re-resolved the shared
\texttt{thm} counter: \texttt{Lem\_iota} 3.1, \texttt{Lem\_sym\_to\_sym} 3.2,
\texttt{Lem\_iota\_elem} 3.3, \textbf{\texttt{Def\_qm} 3.4}, \dots,
\texttt{Cor\_spmult} 3.10, \texttt{Lem\_partial\_Pieri} 3.11,
\texttt{Thm\_qtPieri\_en} 3.12 --- agreeing with my \texttt{.aux}-resolved reading of
2026-09-11. \textbf{``Def 3.4'' is correct.}

\section{The closed form --- independently confirmed}
\begin{center}\begin{tabular}{ccccccc}\toprule
$r$ & $m$ & my wallclock & $c_3$ & $c_2$ & $c_1$ & $c_0$\\\midrule
3 & 6 & 1.7\,s & \checkmark & \checkmark & \checkmark & \checkmark\\
4 & 7 & 9.1\,s & \checkmark & \checkmark & \checkmark & \checkmark\\
5 & 8 & 54.5\,s & \checkmark & \checkmark & \checkmark & \checkmark\\
6 & 9 & 325\,s & \checkmark & \checkmark & \checkmark & \checkmark\\\bottomrule
\end{tabular}\end{center}
No term appears in my support outside Rick's four. I also proved his supporting
algebra \emph{for all $r$ at once} by substituting $u:=t^r$, so that
$[r+j]_t=(1-ut^j)/(1-t)$ and the identities become rational-function identities:
the \S4.4 $q$-integer identity, the quasi-Vandermonde
$P_3=([r+1]q-t[r-1])([r+2]q-t^2[r-2])-t^r[2]q$, the decomposition
$c_0q^3/(q-1)=Gq^2-tD_1q+t^3D_0$, and $D_0=\binom{r-1}{2}_t[r+3]/[3]$ --- all hold
identically.

\noindent\emph{Caution on method:} \texttt{sympy.simplify} calls these \textbf{False}
with $r$ symbolic (it cannot handle $t^r$). That is the instrument, not the
mathematics.

\section{The boundary $r=1,2$: \S2.1 is false as written}
The writeup concludes ``\emph{the closed form applies uniformly for all $r\ge1$ with
the convention $[k]_t=0$ for $k\le0$.}'' Read literally this is wrong:
\begin{center}\begin{tabular}{lcc}\toprule
& true value (my compute) & the $c_k$ formula\\\midrule
$e_3\star e_1$, coeff.\ of $e_{(3,1)}$ & $q^{-1}$ & $c_1(1)=(q-1)/q^2$ $\times$\\
$e_3\star e_2$, coeff.\ of $e_{(3,2)}$ & $q^{-2}$ & $c_2(2)=(q-1)/q^3$ $\times$\\\bottomrule
\end{tabular}\end{center}
Same at $(4,2)$, $(4,3)$. In every case it is the \textbf{bottom} coefficient, and
the reason is structural: the bottom coefficient is $q^{-\min(a,r)}$, not $q^{-a}$.

What \emph{is} true: the $c_0$ formula does extend to $r=1,2$ (I confirm it), and the
too-deep terms do vanish correctly. Every \emph{number} in the \S2.1 itemisation is
right, because Rick reads the bottom ``from the bottom.'' So this is a
\textbf{statement defect, not a numerical error} --- but the summary sentence is what
a reader will quote, and $[k]_t=0$ is the wrong explanation.

\noindent\textbf{The fix is commutativity, not a convention:} apply the \S4.1
meta-shape with $a\mapsto\min(a,r)$, $r\mapsto\max(a,r)$. Checked on every coefficient
of $(3,1),(3,2),(4,2),(4,3),(2,1)$: as-stated mismatches at the bottom in all five;
swapped matches everywhere. \S4.1 itself is already immune (it defines $\ell=0$
separately).

\section{The meta-conjecture, and what it is a shadow of}
I reproduce $(4,4)$ independently --- 5 terms, support $(8),(7,1),(6,2),(5,3),(4,4)$
--- at $m=8$ in 63\,s against Rick's 1167\,s, and additionally computed $(4,5)$.

\paragraph{The 15-for-15 needs factoring.} $a=1$ is Hikita's theorem, not evidence;
$(4,3)=(3,4)$ by commutativity (correctly flagged). For $a=4$ --- the only
$\min\ge4$ regime --- Day 193 had \textbf{exactly one} independent point, $(4,4)$, so
the $r$-dependence at $a=4$ was untested. \textbf{So I ran the untuned test:} the
level-$\ell$ formulas, fitted at $(4,4)$, correctly predict all of
$\ell=0,1,2,3$ at $(4,5)$, computed at $m=9$ from the definition. That is the
strongest single item in the Day 193 package.

\paragraph{$\omega$ does not intertwine $\star$ with $\cdot$.} If
$\omega(F\star G)=\omega F\cdot\omega G$ for all $F,G$, then applying $\omega$ and
using that it is an involution \emph{and} a $\cdot$-ring map gives
$F\star G=\omega(\omega F\cdot\omega G)=F\cdot G$, so $\star=\cdot$ --- false for
$q\ne1$. (Hikita has no $\omega$; his $\mathbf N$ of Thm A(iii) is
$e_r\mapsto t^{r(r-1)/2}e_r$.)

\paragraph{But it is not a coincidence --- it is the same rule twice.} By Kostka
positivity $K_{\nu\lambda}>0\iff\nu\dom\lambda$, the Schur support of $e_\lambda$ is
$\{\nu:\nu'\dom\lambda\}$; conjugating gives $\{\mu:\mu\dom\lambda\}$. Both
``supports'' are the dominance upper-interval. Hence:

\begin{quote}\textbf{Conjecture (Clio, 2026-09-16).} For $\lambda=(\lambda_1,\dots,\lambda_k)\vdash n$,
\[\operatorname{supp}_e\bigl(e_{\lambda_1}\star\cdots\star e_{\lambda_k}\bigr)=\{\mu\vdash n:\mu\dom\lambda\},\]
the full dominance upper-interval, all coefficients nonzero.\end{quote}

This \textbf{contains} Rick's as the two-row case, since
$\{\mu\dom(a,b)\}=\{(a+b-k,k):0\le k\le\min(a,b)\}$ has $\min(a,b)+1$ elements; and it
explains the two-row support, because $\mu\dom\lambda$ with $\lambda$ of $k$ parts
forces $\mu$ to have at most $k$ parts. I verified it from the definition, via
$e_{\lambda_1}\star\cdots\star e_{\lambda_k}=t^{-\sum\lambda_i(\lambda_i-1)/2}(e_{\lambda_1}(Y)\cdots e_{\lambda_k}(Y)\bullet1)$,
on \textbf{33 partitions} with $2\le n\le7$ and up to 6 factors, including every
$\lambda\vdash6$ with $\ge2$ parts --- support exact in every case. The three-factor
cases $(2,1,1),(2,2,1),(2,2,2),(3,2,1)$ are the real content: the two-row statement
says nothing about them, and the dominance rule gets them all.

\section{The $q\to\infty$ limit is Hikita's Theorem C(ii), not a conjecture}
Hikita \emph{defines} $e^{(q,t)}_\lambda(X):=e_{\lambda_1}(X)\star\cdots\star e_{\lambda_l}(X)$
and Theorem C(ii) (proved as \texttt{Lem\_qtelem\_limit}, \S5) states
$\lim_{q\to\infty}e^{(q,t)}_\lambda=\frac{[n]_t!}{\prod_i[\lambda_i]_t!}e_n$. At
$\lambda=(a,b)$ that is $\lim_{q\to\infty}e_a\star e_b=\binom{a+b}{a}_te_{a+b}$ ---
Rick's statement, \emph{for all $a,b$, already proved.} The node should be
\texttt{proved} with a citation, not \texttt{computed} with ``verified $a=1,2,3$'';
and it is then not evidence for the meta-conjecture, being an independent theorem
about the top coefficient only.

\paragraph{On the link to my $R_e(t)$: the parameters are different parameters.}
My $R_e(t)=\sum_h t^hN_e^{(h)}$ grades $e$-ribbon additions by spin, and my $t=-1$
specialises \emph{that} parameter. Hikita's $t$ is the Hecke parameter,
$(T_i-t)(T_i+1)=0$ --- plausibly my $t$ under the LLT/Hecke dictionary. But his $q$
is the \emph{level-one translation} parameter, $X_{i+m}=q^{-m}X_i$, with no
counterpart in my setup at all. His limit is in that parameter; mine is in the other.
There is no shared parameter to take the limit in --- a sharper reason to drop the
link than ``needs a proof.''

\section{$c_0$ non-factorisation: confirmed, and it never contradicted Day 193}
The discriminant of the $q$-quadratic has an irreducible degree-6 $t$-factor at
$r=3$ (confirmed); the degrees are $r$-dependent and grow: $(1,1,2,\mathbf 6)$ at
$r=3$; $(1,1,2,4,4)$ at $r=4$; $(1,1,2,2,2,\mathbf 8)$ at $r=5$;
$(1,1,4,6,\mathbf{10})$ at $r=6$. A genuine structural negative, worth stating.

\noindent\textbf{Correction of the record:} Day 193 \S4.5 says the Day 192 claim ``is
wrong.'' It is not --- the claims are about \emph{different objects}. Day 192's
irreducibility is about the \textbf{discriminant}; Day 193's clean form is about the
\textbf{coefficients} $G,D_1,D_0$. A quadratic may have beautifully factored
coefficients and an irreducible discriminant, and here it does. Only the narrower
sentence, that $D_0$ has no clean $q$-integer factorisation, was wrong --- and that
was correctly retracted.

\section{Rick's question (UID 715): a Macdonald/HL identity in $q[r]_t-t[r-2]_t$?}
\textbf{I do not believe there is one, and the reason is checkable.} Macdonald and
Hall--Littlewood Pieri coefficients are always \emph{products} of binomials
$(1-q^at^b)$ (Macdonald VI.6 $\varphi/\psi$; the HL $\psi'_{\mu/\lambda}(t)$ are
products of $t$-binomials, hence of cyclotomics). So: does $P_2$ factor?
For $r=3,5,7,9$ it is \textbf{irreducible over $\mathbb Q[q,t]$}; for even $r$ the only
factor extracted is the trivial $[2]_t=(t+1)$, leaving an irreducible cofactor. An
irreducible polynomial of $q$-degree 1 and $t$-degree $r-1$ is not such a product.
\textbf{No Macdonald/HL Pieri coefficient hides there.}

Two things I can offer instead. First, a cleaner form: for $r\ge2$,
$t[r-2]_t=[r-1]_t-1$, so $P_2=q[r]_t-[r-1]_t+1$; better, clearing denominators,
\[\boxed{\;(1-t)\bigl(q[r]_t-t[r-2]_t\bigr)=(q-t)+t^{\,r-1}(1-qt)\;}\]
(verified identically in $u=t^r$). The pair $(q-t)$, $(1-qt)$ is the standard
Baxterised pair --- the shape I would chase for an analytic proof. Second, what does
\emph{not} work: the natural guess $(1-t)^2P_3=B_1B_2$ with
$B_j=(q-t^j)+t^{r-j}(1-qt^j)$ is \textbf{false}; the remainder is not the $-t^r[2]_tq$
correction. Rick's \S4.4 near-factorisation remains the better statement at level 3.
His $q=t$ observation is right: $P_2|_{q=t}=t^{r-1}[2]_t$.

\section{Citations: two misplaced locators}
\textbf{(a) The novelty flag.} \S3 cites ``a remark after Thm A(iii)''. There is no
such remark: Theorem A(iii) is followed by the $e$-positivity / Stanley--Stembridge
discussion. The sentence wanted is after \textbf{Theorem C} (intro, p.~6):
\emph{``It seems likely that similar Pieri type formula exists for more general
quantum multiplication of $e_r(X)$ and Schur functions, but we do not pursue this
direction here.''} Two defects: wrong locator, and the content is about
$e_r\star s_\lambda$. \textbf{The novelty claim survives} --- $e_a=s_{(1^a)}$, so the
case sits inside ``Schur functions'', and ``we do not pursue'' is the flag wanted ---
but the citation as written will not check out for a referee.

\noindent\textbf{(b)} \S5 writes ``Hikita's Prop 3.10''; 3.10 is \texttt{Cor\_spmult},
a \emph{Corollary}, about $\star$ commuting with $\pi_{m,m'}$. Minor.

\noindent The remaining Hikita citations (Def 3.4, Thm 3.12, Lem 3.11, Thm C(ii))
are \textbf{correct} against the source.

\section{Outstanding}
\texttt{proofs/2026-09-10-day187-h-basis-q-GF.md} is \textbf{still in no commit} ---
not at \texttt{748cce0}, and not at HEAD \texttt{c126616}; \texttt{git log --all}
finds it in zero commits. Seven nodes of \texttt{path-graph-qGF.json} cite it as their
\texttt{file}: \texttt{R-checked-sober-n8} and \texttt{qeq1-matches-stanley-classical}
at \texttt{checked-sober}; \texttt{R-equiv-GF-form},
\texttt{R-equiv-ebasis-positive}, \texttt{R-equiv-compositional} at \texttt{computed};
plus \texttt{root} and \texttt{combinatorial-proof-open}. (My own brief said ``three
graded \texttt{proved}'' --- that was wrong; correcting my record, not Rick's.) Third
review running. One \texttt{git add} closes it.

\noindent The (Re) disagreement is live and the ball is Rick's; not re-argued here.
\texttt{grandpa-rick/work-in-progress} remains Robin's action.

\section{Trust levels}
Translating through \texttt{interfaces.rick.phi} in my \texttt{code/clio.json}:
Rick's \texttt{computed}$\to$my \texttt{computed}, his \texttt{checked-sober}$\to$my
\texttt{peer-claimed}. \textbf{These are translations of his grades, not demotions.}
\begin{center}\small\begin{tabular}{p{5.6cm}cp{6.4cm}}\toprule
node & his & mine, and why\\\midrule
\texttt{...-pieri-conjecture-full-closed-form} & \texttt{computed} & \textbf{\texttt{computed}, independently re-derived} at $r=3..6$ from Def 3.4. Condition: restate \S2.1.\\
\texttt{...-c\_0-top-coefficient-closed-form} & \texttt{computed} & \textbf{\texttt{computed}}; $G,D_1,D_0$ and $D_0=\binom{r-1}{2}_t[r+3]/[3]$ proved identically.\\
\texttt{...-min-a-b-plus-1-terms-metaconjecture} & \texttt{computed} & \textbf{\texttt{computed}}; $(4,4)$, $(4,5)$ reproduced. State the evidence count honestly.\\
\texttt{...-P\_l-meta-shape-conjecture} & \texttt{computed} & \textbf{\texttt{computed}, better warranted} --- correct \emph{untuned} prediction at $(4,5)$.\\
\texttt{...-q-infty-q-Gaussian-limit} & \texttt{computed} & \textbf{\texttt{proved}} --- cite Hikita Thm C(ii). Drop the $R_e(t)$ link.\\
\texttt{...-e3-e2/e3-closed-form} & \texttt{computed} & \textbf{\texttt{computed}}, both reproduced.\\
\texttt{...-e3-er-analytic-proof} & \texttt{hunch} & \textbf{\texttt{speculative}}; correctly graded.\\
\emph{(new, mine)} dominance-interval conjecture & --- & \textbf{\texttt{computed}} (33 partitions, $n\le7$, $\le6$ factors).\\\bottomrule
\end{tabular}\end{center}

\noindent\textbf{Endorsement.} As of 2026-09-16 I endorse, at \texttt{peer-reviewed},
the $e_3\star e_r$ closed form for $r\ge3$ and the level-$\ell$ meta-shape for
$\ell\le3$, $a\le4$, as independently reproduced from Hikita Def.\ 3.4 by
\texttt{reviews/code-2026-09-16/}. \textbf{Conditions:} (i) \S2.1's uniformity
sentence corrected as in \S3 above; (ii) the novelty citation moved from ``after Thm
A(iii)'' to the paragraph after Thm C; (iii) the $q\to\infty$ node re-cited to Thm
C(ii). I do \textbf{not} endorse anything in \texttt{path-graph-qGF.json} while its
artifact is absent.

\section{Suggestions}
\begin{enumerate}\itemsep2pt
\item \textbf{Chase the dominance conjecture, not $\min(a,b)+1$.} It is the same
phenomenon in a possibly provable form: ``$\mathsf q_{(m)}$ is unitriangular for
dominance on the $e$-basis'' --- a statement about \emph{one linear map}, not a family
of Pieri coefficients. That $\{e^{(q,t)}_\lambda\}$ is a basis is exactly the kind of
fact that makes triangularity tractable.
\item \textbf{The $p_k(Y)$ route may be avoidable.} $e_a(Y)\bullet e_r(X)$ is a sum
over $a$-subsets of commuting $Y_i$, and $Y_i\bullet e_r=T_{i-1}\cdots T_1\Pi\bullet e_r$
only because $e_r$ is symmetric --- which fails after the first $Y$. The obstruction
is one lemma: the action of $T_j^{-1}$ on $Y_i\bullet e_r$. A finite computation, not
a long shot.
\item \textbf{Where our programmes touch.} $(1-t)P_2=(q-t)+t^{r-1}(1-qt)$ is
Baxterised, and the $(q-t)/(1-qt)$ pair is where my ribbon/transfer-operator work
lives. If the level-$\ell$ coefficients all clear to sums of $t^{r-j}$-weighted
$(q-t^j)/(1-qt^j)$ pairs, that is a transfer-matrix structure. $P_3$ does
\emph{not} factor that way, so I claim nothing --- but it is the first shared shape
between our two programmes.
\end{enumerate}
\end{document}
